\section{Architecture and Topology}\label{sec:arch}
% Rubric criterion 2 (15 %): diagrams with logical/physical layers, demarcation points, VLANs and IP boundaries labelled.

The network has three layers of responsibility, and the figures in this section follow them:
\emph{campus} (switching and gateways inside a site), \emph{site routing} (OSPF between the gateways and
the edge routers of one site) and \emph{WAN} (BGP between the sites and towards the providers). The
editable sources of all diagrams are in \path{Design/TVI_Topology.drawio}.

\subsection{Enterprise overview}\label{sec:arch-overview}
\Cref{fig:overview} shows the three sites, the two providers and the overlay. Each site has two edge
routers; edge router~1 of every site connects to ISP-A (AS~\FactAsIspA) and edge router~2 to ISP-B
(AS~\FactAsIspB). The \emph{demarcation point} between TVI and a provider is the /30 on interface
\texttt{Gi0/0} of each edge router: everything to the left of it is TVI's responsibility and addressing,
everything to the right belongs to the provider. Both providers connect to a common Internet router that
stands for the public cloud (SaaS and public DNS).

The overlay consists of \FactNumTunnels{} encrypted tunnels in a hub-and-spoke arrangement. The two HQ edge
routers are the hubs; each branch has one tunnel through ISP-A to HQ edge router~1 and one through ISP-B to HQ
edge router~2. A branch therefore survives the loss of either provider, either of its own edge routers, or
either HQ edge router.

\diagram[fig:overview]{01_wan_overview}{Enterprise overview: three sites, two ISPs, GRE/IPsec overlay and demarcation points}{\textwidth}

\textbf{Why hub-and-spoke and not a full mesh.} The data centre, the DHCP/DNS service and most
applications are at HQ, so almost all inter-site traffic is branch-to-HQ. A full mesh over both providers
would need twelve tunnels instead of four for the rare branch-to-branch flows, which here take one extra
hop through an HQ edge router. Dynamic spoke-to-spoke tunnels (DMVPN) are the natural next step and are
listed in the roadmap; a controller-based SD-WAN would automate the same overlay (\cref{sec:sdwan}).

\subsection{Headquarters campus (HCMC)}\label{sec:arch-hq}
\Cref{fig:hq} shows HQ with interface identifiers. The campus uses a \emph{collapsed core}: two
Layer~3 switches act as distribution and core at once. They are joined by a two-member LACP port-channel
(Po1) that carries all VLANs, and each access switch has one uplink to each core. The \emph{L2/L3
boundary} is at the core switches: VLANs exist only between the access ports and the core SVIs. From
the cores upward every link is a routed /30 -- each core connects to both edge routers, and the two edge
routers are connected to each other. The data-centre service node \dev{HCMC-DC-SRV01} sits in VLAN~100 and
provides DHCP, DNS and NTP for the whole company.

We chose a collapsed core because TVI's HQ access layer is a single building: a separate distribution
layer would put two more devices into every path without isolating anything. The routed uplinks already
separate core and edge, so distribution pairs can be inserted later without renumbering
(\cref{sec:ip-capacity}). Keeping Layer~2 inside the campus also confines spanning tree and broadcast
traffic to a small domain -- the broadcast storms of the legacy flat network cannot cross a routed link.

The direct link between the two edge routers (\texttt{Gi0/3}) is not decoration: without it, a single
branch-provider failure causes a routing loop between a core switch and an edge router
(\cref{sec:underlay,sec:challenges}).

\begin{landscape}
\diagram[fig:hq]{02_hq_detail}{HCMC headquarters: interface IDs, VLANs and the L2/L3 boundary}{0.9\linewidth}
\end{landscape}

\subsection{Branch sites (Hanoi, Da Nang)}\label{sec:arch-branches}
The branches (\cref{fig:branches}) are built from two edge routers and an access switch. The routers are the
default gateways of the branch VLANs: each has an 802.1Q trunk to the switch with one sub-interface per VLAN
(router-on-a-stick), and HSRP provides a virtual gateway address shared by both. The active role is split
per VLAN, so in normal operation both routers -- and both providers -- carry traffic. Each router tracks
its route to HQ; if it loses its tunnel, it hands its VLANs to the other router (\cref{sec:fhrp}). The two
routers form their OSPF adjacency over the management VLAN (99).

A pair of Layer~3 switches per branch would be the next step at the 2030 size; today it would double the
branch device count without adding a function. The single access switch is the remaining single point of
failure for the hosts of a branch in the lab, and we list it as a limitation (\cref{sec:conclusion}).

\diagram[fig:branches]{03_branches}{Branch sites: dual edge routers, router-on-a-stick gateways, local Internet breakout}{\textwidth}

\subsection{Routing architecture}\label{sec:arch-routing}
\Cref{fig:routing} separates the routing layers. Inside each site, OSPF (OSPFv3 for IPv6) connects the
gateways and the edge routers and carries a default route from the edge routers to the core. Between the
sites, iBGP in AS~\FactAsTvi{} runs over the tunnels and carries exactly three IPv4 prefixes, one per site;
the HQ edge routers act as route reflectors. Towards the providers, eBGP accepts only a default route and
advertises nothing. \Cref{sec:underlay,sec:overlay} explain why the work is divided this way.

\diagram[fig:routing]{04_routing_design}{Routing design: OSPF underlay per site, eBGP to the ISPs, iBGP route-reflector overlay over GRE/IPsec}{\textwidth}

\subsection{Device inventory}\label{sec:arch-devices}
\Cref{tab:devices} lists the \FactNumDevices{} network devices of the lab; the cabling is in
\cref{app:wiring}.
% GENERATED by Tools/gen_latex.py - do not edit; change Tools/tvi_model.py and regenerate.
\begin{table}[htbp]
\centering\scriptsize
\caption{Device inventory of the lab}\label{tab:devices}
\begin{tabularx}{\textwidth}{l l l l l X}
\toprule
\textbf{Hostname} & \textbf{Site} & \textbf{Role} & \textbf{Image} & \textbf{Mgmt / Lo0} & \textbf{Function} \\
\midrule
\texttt{HCMC-CORE-SW01} & HCMC & core & IOSvL2 & \texttt{10.1.99.2} & Collapsed core / L3 gateway, STP root for odd VLANs \\
\texttt{HCMC-CORE-SW02} & HCMC & core & IOSvL2 & \texttt{10.1.99.3} & Collapsed core / L3 gateway, STP root for even VLANs \\
\texttt{HCMC-ACC-SW01} & HCMC & access & IOSvL2 & \texttt{10.1.99.11} & Access switch (user floors) \\
\texttt{HCMC-ACC-SW02} & HCMC & access & IOSvL2 & \texttt{10.1.99.12} & Access switch (voice + data-centre services) \\
\texttt{HCMC-EDGE-RTR01} & HCMC & edge & IOSv & \texttt{10.1.255.11} & WAN edge, ISP-A, IPsec hub, NAT, ZBF \\
\texttt{HCMC-EDGE-RTR02} & HCMC & edge & IOSv & \texttt{10.1.255.12} & WAN edge, ISP-B, IPsec hub, NAT, ZBF \\
\texttt{HCMC-DC-SRV01} & HCMC & dcsrv & IOSv & \texttt{10.1.100.10} & Centralised DHCP/DHCPv6/DNS/NTP (router acting as server) \\
\texttt{HANOI-EDGE-RTR01} & HANOI & edge & IOSv & \texttt{10.2.255.1} & Branch gateway (HSRP), ISP-A, IPsec spoke \\
\texttt{HANOI-EDGE-RTR02} & HANOI & edge & IOSv & \texttt{10.2.255.2} & Branch gateway (HSRP), ISP-B, IPsec spoke \\
\texttt{HANOI-RND-SW01} & HANOI & access & IOSvL2 & \texttt{10.2.99.11} & R\&D access switch \\
\texttt{DANANG-EDGE-RTR01} & DANANG & edge & IOSv & \texttt{10.3.255.1} & Branch gateway (HSRP), ISP-A, IPsec spoke \\
\texttt{DANANG-EDGE-RTR02} & DANANG & edge & IOSv & \texttt{10.3.255.2} & Branch gateway (HSRP), ISP-B, IPsec spoke \\
\texttt{DANANG-LOG-SW01} & DANANG & access & IOSvL2 & \texttt{10.3.99.11} & Logistics / IoT access switch \\
\texttt{ISP-A-RTR} & provider & isp & IOSv & \texttt{192.0.2.1} & ISP-A (AS64501) \\
\texttt{ISP-B-RTR} & provider & isp & IOSv & \texttt{192.0.2.2} & ISP-B (AS64502) \\
\texttt{INTERNET-RTR} & provider & internet & IOSv & \texttt{192.0.2.3} & Simulated Internet / SaaS (AS64510) \\
\bottomrule
\end{tabularx}
\end{table}

